\chapter{子词分词}
\label{lab:21_tokenizer}

\noindent\textbf{配套文件：}\path{notebook/21_tokenizer.ipynb}\par
\noindent\textbf{教材关联：}\bookxrefshort{ch:rev-tokenization}。

\section*{实验目标}
验证编码可逆性、合并顺序和特殊词元协议。

\section*{准备及定位}
先阅读本册实验总则，确认Notebook中的依赖与资产单元。原理计算、库接口迁移和真实模型训练可能具有不同资源要求，应分别记录设备、版本及运行范围。

源码定位可从下列函数开始；应结合前置数据与配置单元阅读。
\begin{itemize}
\item \texttt{my\_locate\_project\_directory}
\item \texttt{my\_sha256\_file}
\item \texttt{my\_is\_clean\_record}
\end{itemize}

\section*{操作及观察}
\begin{enumerate}
\item 固定训练文本，跟踪高频字节对合并和词元边界变化。
\item 用已见及未见文本检查编码—解码往返；与tiktoken对照时记录词表与预处理差异。
\item 检查BOS、EOS、PAD及聊天或翻译模板，构造批次并区分padding与截断。
\item 保存并重载分词制品，复查特殊ID、样本编码和标签屏蔽。
\end{enumerate}

\section*{结果判读及提交}
提交合并轨迹、往返失败样本、特殊ID表和重载结果。不同训练词表不要求产生相同ID；目标文本的padding不能计入有效损失。

提交时附上所执行单元范围、环境与制品身份、原始结果及失败说明。将未运行项目明确列出，不能用Notebook中已有输出替代本次执行证据。

相关方法的原始定义与适用前提见\cite{sennrichSubwordUnits2016}；本实验的运行结果仍须按实际配置单独验证。
